\documentclass[11pt]{article}
\usepackage[hyperref]{acl}
\usepackage{times}
\usepackage{latexsym}
\usepackage{microtype}
\usepackage{listings}
\usepackage{xcolor}
\lstset{
  basicstyle=\ttfamily\scriptsize,
  breaklines=true,
  breakatwhitespace=false,
  frame=single,
  backgroundcolor=\color{gray!8},
  aboveskip=4pt,
  belowskip=4pt,
  columns=flexible,
}

\title{UnMask: Supplementary Material\\
\large Conversation Transcripts (Appendix A) and VLM Session (Appendix C)}

\author{Sanika Vilas Najan \quad Vaishak Girish Kumar \\
  University at Buffalo \\
  CSE 635: NLP and Text Mining --- Spring 2026}

\begin{document}
\maketitle

\section*{Appendix A: Conversation Transcripts}
\label{app:transcripts}

The P3 specification requires 5 conversation transcripts demonstrating no
answer leak before the 3rd turn, plus multimodal VLM demonstration (Appendix C).
All transcripts were generated live from UnMask with cold-start mastery = 0.20
(PCR \texttt{context\_only} mode).
Gold answers are shown for grader reference only; they were present in
\texttt{InternalAnalysis} (hidden) and absent from \texttt{VisibleResponse}
(shown to student).
Leak status verified by both layers of the answer-leak detector.

\subsection*{Transcript 1 --- Brachial Plexus Origins}

\noindent\textbf{Turn 1} $|$ PCR: \texttt{context\_only} $|$ Leak: None $|$ Ends `?': Yes\\
\textit{Student:} What spinal cord levels form the brachial plexus?\\
\textit{UnMask:} Great question! Let's explore the spinal levels involved.
Which specific cervical and thoracic spinal nerves contribute to the formation
of the brachial plexus, and how do they relate to the vertebrae they emerge
from?\\
{\small\textit{[Gold held internally: ``C5--T1 anterior rami'']}}

\noindent\textbf{Turn 2} $|$ PCR: \texttt{context\_only} $|$ Leak: None $|$ Ends `?': Yes\\
\textit{Student:} Is it only cervical nerves?\\
\textit{UnMask:} You're on the right track! Have you considered whether any
thoracic spinal nerves might also contribute to the brachial plexus, in
addition to the cervical nerves?\\
{\small\textit{[Gold held internally: ``Lower 4 cervical (C5--C8) and first thoracic (T1)'']}}

\noindent\textbf{Turn 3} $|$ PCR: \texttt{prerequisite\_first} $|$ Leak: None $|$ Ends `?': Yes\\
\textit{Student:} So five roots total from C5 to T1?\\
\textit{UnMask:} You're doing well in identifying the spinal roots involved!
How do these spinal nerve roots organize themselves into the trunks and
divisions of the brachial plexus?\\
{\small\textit{[Gold confirmed: ``Yes---C5, C6, C7, C8, T1'']}}

\subsection*{Transcript 2 --- Axillary Nerve \& Shoulder Dislocation}

\noindent\textbf{Turn 1} $|$ PCR: \texttt{context\_only} $|$ Leak: None $|$ Ends `?': Yes\\
\textit{Student:} Which nerve is damaged when a patient can't abduct past 15 degrees after shoulder dislocation?\\
\textit{UnMask:} What is the primary muscle responsible for shoulder abduction beyond 15 degrees, and which nerve innervates this muscle?\\
{\small\textit{[Gold held: ``Axillary nerve'']}}

\noindent\textbf{Turn 2} $|$ PCR: \texttt{context\_only} $|$ Leak: None $|$ Ends `?': Yes\\
\textit{Student:} So is it the axillary nerve?\\
\textit{UnMask:} You're on the right track! What is the primary muscle responsible for shoulder abduction beyond 15 degrees, and which nerve innervates this muscle?\\
{\small\textit{[Gold held: ``Axillary nerve innervates deltoid, crucial for abduction beyond 15°'']}}

\noindent\textbf{Turn 3} $|$ PCR: \texttt{prerequisite\_first} $|$ Leak: None $|$ Ends `?': Yes\\
\textit{Student:} What else does the axillary nerve innervate besides deltoid?\\
\textit{UnMask:} Besides the deltoid, which other muscle and what specific area of skin does the axillary nerve innervate?\\
{\small\textit{[Gold confirmed: ``Teres minor + lateral shoulder skin (regimental badge area)'']}}

\subsection*{Transcript 3 --- Radial Nerve \& Wrist Drop}

\noindent\textbf{Turn 1} $|$ PCR: \texttt{context\_only} $|$ Leak: None $|$ Ends `?': Yes\\
\textit{Student:} A patient has wrist drop after a mid-shaft humerus fracture. What nerve is likely injured?\\
\textit{UnMask:} What is the relationship between the location of a mid-shaft humerus fracture and the nerve that runs in that area, particularly in relation to wrist extensor function?\\
{\small\textit{[Gold held: ``Radial nerve'']}}

\noindent\textbf{Turn 2} $|$ PCR: \texttt{context\_only} $|$ Leak: None $|$ Ends `?': Yes\\
\textit{Student:} Does the radial nerve innervate all wrist extensors?\\
\textit{UnMask:} What is the role of the radial nerve in the forearm, and how does it relate to the muscles responsible for wrist extension?\\
{\small\textit{[Gold held: ``Radial nerve innervates all wrist extensor muscles'']}}

\noindent\textbf{Turn 3} $|$ PCR: \texttt{prerequisite\_first} $|$ Leak: None $|$ Ends `?': Yes\\
\textit{Student:} What about sensation: where would the patient feel numbness?\\
\textit{UnMask:} Which specific area on the dorsal side of the hand and forearm would you expect to experience numbness if the radial nerve is injured?\\
{\small\textit{[Gold confirmed: ``Posterior forearm and dorsal hand; first dorsal web space'']}}

\subsection*{Transcript 4 --- Rotator Cuff Anatomy}

\noindent\textbf{Turn 1} $|$ PCR: \texttt{context\_only} $|$ Leak: None $|$ Ends `?': Yes\\
\textit{Student:} Name the four rotator cuff muscles.\\
\textit{UnMask:} What acronym can help you remember the names of the four muscles that form the rotator cuff, and what does each letter stand for?\\
{\small\textit{[Gold held: ``Supraspinatus, infraspinatus, teres minor, subscapularis (SITS)'']}}

\noindent\textbf{Turn 2} $|$ PCR: \texttt{context\_only} $|$ Leak: None $|$ Ends `?': Yes\\
\textit{Student:} I know supraspinatus, but what about the others on the posterior scapula?\\
\textit{UnMask:} What are the names of the rotator cuff muscles that originate from the infraspinous fossa and the teres minor region on the posterior scapula?\\
{\small\textit{[Gold held: ``Infraspinatus and teres minor'']}}

\noindent\textbf{Turn 3} $|$ PCR: \texttt{prerequisite\_first} $|$ Leak: None $|$ Ends `?': Yes\\
\textit{Student:} What movement does the infraspinatus produce?\\
\textit{UnMask:} Considering the infraspinatus is located behind the glenohumeral joint, how might its position influence the type of rotation it facilitates at the shoulder?\\
{\small\textit{[Gold confirmed: ``External rotation of the shoulder'']}}

\subsection*{Transcript 5 --- Median Nerve \& Carpal Tunnel Syndrome}

\noindent\textbf{Turn 1} $|$ PCR: \texttt{context\_only} $|$ Leak: None $|$ Ends `?': Yes\\
\textit{Student:} What nerve is compressed in carpal tunnel syndrome?\\
\textit{UnMask:} Which nerve passes through the carpal tunnel along with the flexor tendons and is responsible for sensation in the thumb, index, middle, and part of the ring finger?\\
{\small\textit{[Gold held: ``Median nerve'']}}

\noindent\textbf{Turn 2} $|$ PCR: \texttt{context\_only} $|$ Leak: None $|$ Ends `?': Yes\\
\textit{Student:} What symptoms would the patient have?\\
\textit{UnMask:} When the median nerve is compressed at the wrist, what changes might you expect in both sensation and muscle function in the hand?\\
{\small\textit{[Gold held: ``Numbness/tingling digits 1--3.5; thenar weakness'']}}

\noindent\textbf{Turn 3} $|$ PCR: \texttt{prerequisite\_first} $|$ Leak: None $|$ Ends `?': Yes\\
\textit{Student:} Would the patient have weakness gripping or just sensory loss?\\
\textit{UnMask:} What role do the thenar muscles, innervated by the median nerve, play in gripping, and how might their dysfunction affect hand strength?\\
{\small\textit{[Gold confirmed: ``Both---thenar weakness impairs pinch/grip; sensory loss in lateral digits'']}}

\bigskip
\noindent\textbf{Transcript summary:} 15 turns, 5 transcripts, 0 confirmed leaks,
0 soft flags, 15/15 ends with ``?''.
Gold answer present only in \texttt{InternalAnalysis} (stripped before
rendering); absent from retrieved Qdrant chunks due to PCR \texttt{must\_not}
filter.

%% ─────────────────────────────────────────────────────────────
\section*{Appendix B: Onboarding, Visual Aid, and Session Summary}

\subsection*{B.1 Personalized Onboarding and Diagnostic Reordering}

The welcome prompt asks a single conversational question; the student's reply
is parsed for \texttt{study\_focus} and \texttt{learning\_mode} before the
LangGraph graph is invoked.

\noindent\textbf{UnMask welcome:}\\
\textit{``Hey! Welcome to UnMask. I'm your NBCOT anatomy study companion. I
use the Socratic method to help you build clinical reasoning, not just recall
facts. What are you finding toughest right now, and do you prefer learning
with diagrams or written Q\&A? (Just reply naturally, I'll calibrate from
there!)''}

\noindent\textbf{Student:} \textit{``I keep getting confused by the brachial
plexus, especially which levels form it. I think I learn better with
diagrams.''}

\noindent\textbf{System action} (Orchestrator, before \texttt{graph.invoke}):
\begin{lstlisting}
study_focus = "brachial plexus"
learning_mode = "visual"
diag_order = [0, 2, 3, 1]  # BP questions first
visual_hint_threshold = 1   # diagram on 1st wrong
\end{lstlisting}

\noindent\textbf{Diagnostic Q1 (reordered)} $|$ Phase: \texttt{rapport}\\
\textit{UnMask:} ``Quick diagnostic -- 4 questions to calibrate where we start:
\textbf{Q1 of 4:} Which spinal cord levels give rise to the brachial plexus?''\\
\textit{Student:} ``C5 to T1'' $\to$ \textbf{Correct.}
\texttt{mastery[brachial\_plexus.origin] = 0.50}

\subsection*{B.2 Visual Aid Trigger (Gray's Anatomy Plate)}

Visual mode threshold is 1. Below: first incorrect answer in Tutoring phase.

\noindent\textbf{Turn 1} $|$ PCR: \texttt{context\_only} $|$ Phase: \texttt{tutoring}\\
\textit{UnMask:} ``How do the C5--T1 roots combine to form the three trunks
of the brachial plexus?''\\
\textit{Student:} ``I'm not sure.'' $\to$ \textbf{Incorrect.}
\begin{lstlisting}
consecutive_incorrect = 1  # >= threshold (1)
\end{lstlisting}

\noindent\textbf{System action} (animated HTML diagram, served via FastAPI \texttt{/static/anatomy/}):
\begin{lstlisting}
# api.py emits SSE event:
{"type": "visual_hint",
 "concept": "brachial_plexus",
 "image_url": "/static/anatomy/brachial_plexus.html",
 "caption": "Brachial Plexus (C5-T1)"}
# Frontend DiagramCard renders as animated <iframe>
\end{lstlisting}

\noindent\textbf{UnMask reply with inline plate:}\\
\textit{``That's a tricky one to visualize! Here is the brachial plexus
diagram for reference. [Gray809 image renders inline.] Look at where C5--C6
converge: what structure do they form at that junction?''}

\subsection*{B.3 Session Wrapup Summary (Structured Output)}

At wrapup ($t \geq 840$s), GPT-4o structured output produces
a \texttt{SessionSummary}:

\begin{lstlisting}
overall_assessment:
  "Solid C5-T1 recall and correct radial wrist drop.
   Trunk formation and teres minor need more work."

topic_reports:
  - concept: brachial_plexus.trunks
    mastery_score: 0.25
    status: needs_review
    honest_feedback: "Struggled turns 1+4; review
      how C5-C6 form the upper trunk."
  - concept: brachial_plexus.origin
    mastery_score: 0.50
    status: progressing
    honest_feedback: "Root levels correct; reinforce
      anterior rami distinction."
  - concept: peripheral_nerves.radial
    mastery_score: 0.72
    status: mastered
    honest_feedback: "Wrist drop identified correctly
      with good clinical reasoning."

mistake_highlights:
  - "Confused upper trunk (C5+C6) with C5 alone."
  - "Missed teres minor (2nd axillary target)."

study_recommendations:
  - "Draw RTTDC (Roots-Trunks-Divisions-Cords)."
  - "Practice Phalen's and Tinel's sign sequences."

youtube_resources:
  - {concept: "brachial_plexus.trunks",
     title: "Brachial Plexus in 5 Minutes",
     creator: "Ninja Nerd Science",
     search_query: "brachial plexus trunks anatomy NBCOT"}
  - {concept: "brachial_plexus.origin",
     title: "Spinal Nerves and the Brachial Plexus",
     creator: "Armando Hasudungan",
     search_query: "spinal nerve roots brachial plexus C5 T1"}

closing_reflection:
  "Patient has deltoid weakness and lateral arm
   numbness post-dislocation. Which nerve and cord
   would you evaluate first?"
\end{lstlisting}

\noindent The \texttt{mistake\_highlights} field is populated from
\texttt{mistake\_log} entries whose \texttt{misconception} strings were
captured at the time of each incorrect answer (hidden
\texttt{InternalAnalysis.student\_misconception}).
The \texttt{youtube\_resources} field contains 2--4 resources for the weakest
concepts, rendered as clickable YouTube search cards in the frontend Progress tab.
The \texttt{closing\_reflection} ends with ``?'' --- the session ends on a
Socratic note, not a declarative summary.

%% ─────────────────────────────────────────────────────────────
\section*{Appendix C: Multimodal VLM Session (Task 4)}
\label{app:vlm}

\noindent\textbf{Context:} Student uploads a photo of a labeled anatomical diagram
of the brachial plexus via the Next.js frontend file picker during the tutoring
phase (consecutive\_incorrect $= 1$, PCR mode: \texttt{context\_only}).
The image is base64-encoded by the frontend and POSTed to
\texttt{/api/sessions/\{id\}/upload}; the backend calls \texttt{analyze\_uploaded\_image}.

\noindent\textbf{System action} (\texttt{analyze\_uploaded\_image}, Gemini 2.0 Flash Lite via OpenRouter):
\begin{lstlisting}
Image received: student_upload_brachial.jpg (base64)
VLM internal: "Identified structure: Brachial plexus diagram
  showing C5-T1 roots, three trunks (upper/middle/lower),
  and lateral/medial/posterior cords."
VLM output (Socratic, structure unnamed):
  "I can see a nerve network diagram here --- look at the
   connections between the spinal levels at the top and
   the groupings below them. What do you notice about
   how the upper levels combine compared to the lower
   ones, and what might that tell you about how injury
   at different levels would produce different symptoms?"
\end{lstlisting}

\noindent\textbf{Turn 1} $|$ Source: student image upload $|$ VLM: Gemini 2.0 Flash Lite $|$ Leak: None\\
\textit{Student:} [uploads anatomical image]\\
\textit{UnMask (VLM):} ``I can see a nerve network diagram here --- look at the connections between the spinal levels at the top and the groupings below them. What do you notice about how the upper levels combine compared to the lower ones, and what might that tell you about how injury at different levels would produce different symptoms?''\\
{\small\textit{[Structure identified internally: brachial plexus roots/trunks --- not named in response]}}

\noindent\textbf{Turn 2} $|$ Phase: \texttt{tutoring} $|$ PCR: \texttt{context\_only} $|$ Leak: None\\
\textit{Student:} ``The upper levels C5 and C6 seem to merge together.''\\
\textit{UnMask:} ``Exactly right --- they do converge. Now, if that convergence point were injured, which specific movements of the arm would you expect to be affected, and why?''\\
{\small\textit{[Gold held: ``C5+C6 form upper trunk; injury causes Erb's palsy --- shoulder abduction/external rotation lost'']}}

\noindent\textbf{Turn 3} $|$ PCR: \texttt{prerequisite\_first} $|$ Leak: None\\
\textit{Student:} ``Shoulder movement would be lost?''\\
\textit{UnMask:} ``You're on the right track! Specifically, which shoulder movements depend on C5 and C6, and what clinical presentation would you expect in a newborn with a difficult delivery?''\\
{\small\textit{[Gold confirmed: ``Waiter's tip posture (Erb's palsy) --- abduction and external rotation lost'']}}

\bigskip
\noindent\textbf{VLM performance note:} The structure was identified without being named in the Socratic question (0 leaks). The VLM response maintained Socratic form (ends with ``?'') and guided the student toward discovering the convergence pattern independently. This completes Task 4 (Multimodal Diagram Tutoring) from the project specification.

\end{document}
